\chapter[\emj{🎄}~Segment Tree (Árbol de Segmentos)]{\emj{🎄}~Segment Tree (Árbol de Segmentos)}

\begin{dsaquees}

Un array donde te preguntan miles de veces "¿cuál es la suma (o el
mínimo) entre la posición 3 y la 8?" y ADEMÁS cambian valores entre
pregunta y pregunta 🔄.

Recalcular sumando a mano cada vez cuesta $O(n)$ por pregunta: lento. El
Segment Tree precalcula las respuestas de bloques y las organiza en un
árbol. Así responde cualquier rango en $O(\log n)$ y también actualiza en
$O(\log n)$. Cada nodo del árbol "resume" un pedazo del array; la raíz
resume todo, y las hojas son los elementos individuales.

\end{dsaquees}

\begin{dsaotraforma}

Un Segment Tree es un árbol binario donde cada nodo guarda información
agregada (suma, mínimo, máximo, gcd\ldots) de un segmento contiguo del
array. La raíz cubre $[0, n-1]$; cada nodo se divide en dos mitades hasta
llegar a hojas de un solo elemento.

Permite dos operaciones en $O(\log n)$:
\begin{itemize}
\item \textbf{Query de rango}: combinar los nodos que cubren $[l, r]$.
\item \textbf{Update de punto}: cambiar un elemento y propagar hacia la
      raíz.
\end{itemize}

Con \textbf{lazy propagation} soporta también updates de RANGO (sumar a
todo un intervalo) en $O(\log n)$. Se almacena en un array de tamaño
$\sim 4n$.

\end{dsaotraforma}

\begin{dsadiagrama}
\begin{verbatim}
Array: [2, 1, 5, 3]   (Segment Tree de SUMA)

              [0..3]=11
             /          \
       [0..1]=3        [2..3]=8
        /    \          /    \
   [0]=2  [1]=1    [2]=5   [3]=3

Query suma(1..2) = combina [1]=1 y [2]=5 = 6
Update pos 3 a 10: actualiza hoja [3] y sube: [2..3]=15, raíz=18
Cada consulta/actualización toca O(log n) nodos.
\end{verbatim}
\end{dsadiagrama}

\begin{lstlisting}[language=C++]
CÓMO FUNCIONA (suma, array de tamaño 4n)
- build(node, l, r): si l==r, hoja = arr[l]; si no, construye hijos y
  tree[node] = tree[izq] + tree[der].
- query(node, l, r, ql, qr): si el rango no toca -> 0; si está contenido
  -> tree[node]; si no, combina las dos mitades.
- update(node, l, r, pos, val): baja a la hoja, cámbiala y recombina.

📝 PSEUDOCÓDIGO (query de rango)
──────────────────────────────────────────────────────────────

función query(nodo, l, r, ql, qr):
    SI qr < l O r < ql: return 0            // fuera de rango
    SI ql <= l Y r <= qr: return tree[nodo] // totalmente dentro
    mid = (l + r) / 2
    return query(2*nodo,   l,     mid, ql, qr)
         + query(2*nodo+1, mid+1, r,   ql, qr)

💻 CÓDIGO C++ (Segment Tree de suma)
──────────────────────────────────────────────────────────────

#include <vector>
using namespace std;

struct SegTree {
    int n; vector<long long> tree;
    SegTree(vector<int>& a) { n = a.size(); tree.assign(4*n, 0); build(a,1,0,n-1); }

    void build(vector<int>& a, int node, int l, int r) {
        if (l == r) { tree[node] = a[l]; return; }
        int mid = (l + r) / 2;
        build(a, 2*node, l, mid);
        build(a, 2*node+1, mid+1, r);
        tree[node] = tree[2*node] + tree[2*node+1];
    }
    // suma en [ql, qr]
    long long query(int node, int l, int r, int ql, int qr) {
        if (qr < l || r < ql) return 0;               // sin solape
        if (ql <= l && r <= qr) return tree[node];     // contenido
        int mid = (l + r) / 2;
        return query(2*node, l, mid, ql, qr)
             + query(2*node+1, mid+1, r, ql, qr);
    }
    // arr[pos] = val
    void update(int node, int l, int r, int pos, int val) {
        if (l == r) { tree[node] = val; return; }
        int mid = (l + r) / 2;
        if (pos <= mid) update(2*node, l, mid, pos, val);
        else            update(2*node+1, mid+1, r, pos, val);
        tree[node] = tree[2*node] + tree[2*node+1];
    }
};

⏱️ COMPLEJIDAD (Big-O)
──────────────────────────────────────────────────────────────

  Operación       │ Tiempo     │ Espacio
  ────────────────┼────────────┼──────────
  Build           │ O(n)       │ O(4n)
  Query de rango  │ O(log n)   │ O(log n) pila
  Update de punto │ O(log n)   │ O(log n) pila

* Con lazy propagation, el update de RANGO también es O(log n).
\end{lstlisting}

\begin{dsaentrevistas}

✅ ¿Cuándo usarlo? Muchas consultas de rango CON actualizaciones
   intercaladas. Si el array NUNCA cambia, un prefix-sum simple ($O(1)$
   por query) es mejor y más fácil.

💡 "Range Sum Query - Mutable" (LeetCode 307): el caso canónico.
   Actualizaciones + sumas de rango = Segment Tree o Fenwick.

⚠️ Tamaño 4n: reserva $4n$ para el árbol, no $2n$; con $n$ que no es
   potencia de 2 puede desbordar si escatimas memoria.

🔑 Segment Tree vs Fenwick: Fenwick (BIT) es más corto y rápido para
   sumas de prefijo, pero el Segment Tree es más general (mínimo, máximo,
   gcd, y updates de rango con lazy).

\end{dsaentrevistas}

\begin{dsaresumen}

• Árbol binario; cada nodo resume un segmento del array.
• Query de rango y update de punto en O(log n).
• Raíz cubre todo; hojas son elementos individuales.
• Lazy propagation habilita updates de RANGO en O(log n).
• Se guarda en un array de tamaño 4n.
• Si el array no cambia, usa prefix-sum (más simple).
════════════════════════════════════════════════════════════════

\end{dsaresumen}
